\section{不同模型的最佳实践}
\label{sec:model}

这是本问的核心。要回答「不同模型该怎么做」，先要说清差异从哪来。

\subsection{差异的三个来源}

\begin{table}[t]
\centering
\caption{决定分离做法的三个模型侧变量。三者的乘积决定 KV 传输量与
所需配比，故最佳实践必须按组合分别陈述。}
\label{tab:model-vars}
\small
\begin{tabular}{@{}p{1.35cm}p{2.5cm}p{2.35cm}@{}}
\toprule
变量 & 影响 & 后果 \\
\midrule
KV 每 token 大小 & MLA $<$ GQA $<$ MHA & 决定传输字节数 \\
上下文长度 & 线性放大 KV & 决定是否放得下 \\
负载形态 & 输入/输出长度比 & 决定 P:D 配比 \\
\bottomrule
\end{tabular}
\end{table}

\noindent\textbf{① KV 每 token 大小。} 注意力结构直接决定每个 token 要存多少 KV。
MLA（DeepSeek 系）把 KV 压到低秩潜向量，每 token 的 KV 极小；
GQA（Llama、Qwen 系）用分组共享降到 1/4 至 1/8；MHA 最大。
这个量的差异能达到一个数量级，直接改变传输成本的量级。

实测印证：DeepSeek-V3 的 KV cache 是 250\,MB 量级，经 NVLink 传输
0.6\,ms\texttt{[c-transfer]}。相比之下，同样长度的 MHA 模型 KV 会大出
数倍，传输耗时同步放大。

\noindent\textbf{② 上下文长度。} LoongServe 指出，LLM 的上下文窗口在快速
增长，导致服务系统面临极大的方差，显存消耗随上下文线性上涨
\texttt{[c-model]}。长上下文场景下，KV 不只是「传输成本」问题，
还是「放不放得下」问题——此时分离可以顺带解决容量问题，
而不只是吞吐问题。

\noindent\textbf{③ 负载形态。} 输入长、输出短（如文档问答、代码审查）
意味着 prefill 占主导；输入短、输出长（如创作、推理链）意味着
decode 占主导。前者需要更多 P 卡，后者需要更多 D 卡。
这个比值不是模型属性，是业务属性——同一个模型在不同业务下配比不同。

\subsection{按模型类型的差异化做法}

\noindent\textbf{MLA 系（DeepSeek-V3 等）。} KV 极小，传输成本低，
故分离的门槛低、同机分离几乎白赚。但要注意：正因为 KV 小，
\textbf{分离的收益也随之变小}——收益侧来自配置差异，而 KV 小意味着
decode 的带宽压力相对轻，两阶段配置差异被压缩。
这是一处容易被忽略的反向效应。

\noindent\textbf{MoE 系。} 多了一层专家并行的复杂度。ExpertPlex 报告，
专家共享消除了超过 95\% 的重复模型权重，同时按需复用动态稀疏计算，
而注意力分离则降低了注意力的通信开销\texttt{[c-model]}。
这说明 MoE 场景下的分离不只是 P/D 两分，还要考虑专家在两侧的放置。

\noindent\textbf{稀疏注意力系。} SAC 在 DeepSeek-V3.2 上用 SGLang 做了评测，
报告相对 RDMA 基线取得 2.1× 吞吐、9.7× 更低的 TTFT、1.8× 更低的 TBT，
并据此认为基于 CXL 的分离是稀疏注意力模型更优的基础设施
\texttt{[c-model]}。这条结果的含义是：\textbf{当模型本身的注意力变稀疏，
瓶颈会转移到「KV 放哪里」而非「算得多快」}，此时分离的载体选择
（CXL vs RDMA）比分离本身更重要。

\noindent\textbf{多模态系。} 这是唯一需要把分离从两段变三段的情形。
EAServe 指出，静态 50/50 切分虽然优于不切分，但仍然浪费了本地
prefill 容量\texttt{[c-model]}。原因是多模态有一个 Encode 阶段
（视觉编码），它的资源画像又不同于 prefill 与 decode，
故 SGLang 提供了 EPD（Encode-Prefill-Decode）三段分离。

\noindent\textbf{逐阶段量化。} 一个跨模型通用的手段：既然两阶段偏好不同，
连量化精度也可以不同。但要小心代价——有工作报告，在 prefill 密集型任务上，
只对 prefill 做量化所付出的精度损失，是只对 decode 量化的
1.1 至 4.1 倍（八个模型里有七个如此）\texttt{[c-model]}。
这说明「按阶段分别量化」不是免费的，prefill 侧对精度更敏感。

\subsection{为什么不存在统一配置}

把上面几节合起来看：最优做法取决于 KV 大小（模型）× 上下文长度（业务）
× 输入输出比（业务）× 互联拓扑（部署）× SLO（业务）。五个变量里
只有第一个是模型属性。

因此「某模型的最佳 PD 配置」这个提法本身是有问题的——它缺了四个自变量。
可迁移的结论只能是\textbf{判据}：在给定这五个变量的取值下，用
第 \ref{sec:boundary} 节的条件判断是否分离、按什么比例分离。
